\documentclass[11pt]{article}
\usepackage{fontspec}
\usepackage{amsmath,amssymb,amsthm,mathtools}
\usepackage[margin=1.1in]{geometry}
\usepackage{microtype}
\usepackage{parskip}
\usepackage{lmodern}
\usepackage{xcolor}
\usepackage{adjustbox}
\usepackage[colorlinks=true,linkcolor=blue!50!black,urlcolor=blue!50!black]{hyperref}
\providecommand{\tightlist}{}
\setlength{\emergencystretch}{3.5em}
\title{Causal tomography: certifying manifoldlikeness of a finite causal order from two-dimensional shadows}
\author{Felix Robles Elvira (ORCID: 0009-0009-2017-4394; independent researcher)}
\date{\small preprint, not peer reviewed, version 2026-07-02. This paper is self-contained. The theorems (Sections 2–4) are elementary and proved in-paper given the named imports (Gallai's transitive-orientation theorem, Logan's finite-angle uncertainty principle); the numerical demonstrations are float64 measurement landscapes with fixed seed \texttt{default\_rng(20260702)} and measured-and-disclosed tolerances, and the repository ships the complete demonstration code and canonical outputs (receipts \texttt{i1}–\texttt{i3}, \texttt{k1}–\texttt{k3}, \texttt{l1}–\texttt{l3}, \texttt{m1} in \texttt{code/}, outputs in \texttt{receipts/}). Claims are graded in place: \textbf{[THEOREM]} = proved here; \textbf{[DEMONSTRATED]} = measured with shipped code at the stated scale; verdicts that fall short of a gate are labeled \textbf{inconclusive}, not certified.}
\begin{document}
\maketitle

\section*{Abstract}
Whether a finite partial order is "manifoldlike" — approximable by a uniform-density sprinkling of a Lorentzian spacetime — is the causal-set form of the Hauptvermutung, and above \texttt{1+1} dimensions it collides with an order-theoretic wall: Minkowski dimension detaches from Dushnik–Miller order dimension, and recognition of the relevant order classes is believed hard. We do not attack the decision problem. We split \textbf{verification from search}: we construct a polynomial-time-verifiable \emph{certificate} of higher-dimensional volume-faithfulness — possessing one is proof, up to an exactly characterized budget — and treat finding one as heuristic search whose hardness is expected and priced. Four results. \textbf{(1) A 2D engine [THEOREM].} For orders of dimension ≤ 2 we prove a per-order, deterministic, quantitative characterization: \texttt{P} admits a volume-faithful embedding into the \texttt{M²} diamond with anchored-box defect ≤ ε \textbf{iff} \texttt{dim(P) ≤ 2} and the \emph{canonical rank embedding} — built order-intrinsically from the essentially unique realizer (Gallai) — has star discrepancy \texttt{D* ≤ 3ε + 1/n} (backward direction exact). A relabeling lemma shows the residual realizer freedom (modules) moves the embedded point set not at all; the one exception class (internal prime flips) is closed by census and priced below \texttt{10⁻⁵} where it occurs. \textbf{(2) Causal tomography [THEOREM + DEMONSTRATED].} The order of \texttt{M\textasciicircum{}{1+d−1}} is the intersection of its null-coordinate total orders over the celestial sphere; each antipodal direction pair projects the order to a 2D \textbf{shadow} whose target density is computable in closed form (the chord law \texttt{√((1−|t|)²−s²)} in \texttt{2+1}; the disk-section law \texttt{(1−|t|)²−s²} in \texttt{3+1}). Scoring shadows with a weighted discrepancy against the target copula, plus a \textbf{harmonic condition} — direction profiles of a faithful order are degree-≤1 (Fourier / spherical) harmonics whose coefficients \emph{are} the coordinates — yields certificate conditions that reconstruct coordinates at correlation \texttt{0.9994} per axis in \texttt{2+1} and \texttt{0.9999} per axis in \texttt{3+1} (oracle-family), including, to our knowledge, the first reconstruction of the \textbf{celestial sphere of a causal set from inversion counts alone} (\texttt{S¹}: radius CV \texttt{0.009}; \texttt{S²}: radius CV \texttt{0.008}, angle correlation \texttt{0.9995}). \textbf{(3) The ghost budget [THEOREM-shaped + DEMONSTRATED].} The certificate's soundness is scoped exactly: every marginal-preserving rearrangement is either negligible or a spectral density modulation; a mode invisible to \texttt{K} directions must have frequency \texttt{q ≳ 1/sin(π/2K)} (Logan), and any such mode moves every anchored-box count by at most \texttt{∼a/q} — so the certificate controls the discrepancy metric at all scales, while high-frequency spectral non-uniformity below that amplitude is provably outside \emph{any} finite-\texttt{K} discrepancy certificate. Both halves are demonstrated: a constructed \texttt{q = 120} ghost (genuinely non-uniform, matched-filter \texttt{z = 15.3}) passes the full certificate; the same amplitude at \texttt{q = 6} is caught at over twice the calibrated floor. \textbf{(4) An intrinsic finder [DEMONSTRATED].} A fully order-intrinsic search pipeline — a past/future-count temporal seed, a chain-fleet transverse embedding, an order-only scale calibration, and a \emph{tomographic bootstrap} that uses the reconstruction operator as its own polish — collapses the search gap from \texttt{6×} the certificate floor (annealing from cold) to \texttt{1.5×} at \texttt{N = 400} and to \texttt{1.21×} at \texttt{N = 1000} in \texttt{2+1} — the latter \emph{inside} the certificate gate for the shadow condition — and runs end-to-end in \texttt{3+1} at \texttt{1.50×} (verdict: inconclusive, at exact parity with the \texttt{2+1} pilot's first stage). The same pipeline \emph{refuses} two non-manifoldlike adversary classes. Verification is solved at demonstration scale with its soundness budget stated exactly; search is measurably close; certified end-to-end \texttt{3+1} reconstruction is the stated open problem.

\section*{1. The problem and the reframe}
A causal set is a locally finite partial order proposed as the deep structure of spacetime [1]. Its central mathematical problem is manifoldlikeness: which orders arise (with high probability, or approximately) as Poisson sprinklings of a Lorentzian manifold, with the order given by the causal relation [2, 3]? By Malament's theorem the causal order of a distinguishing spacetime determines its conformal geometry [4], and sprinkling density supplies the volume factor — "order + number = geometry". The verification form of the question, per order: \emph{given one finite order \texttt{P}, decide — and quantify — whether it embeds in a spacetime region at uniform density.}

In \texttt{1+1} dimensions the order theory cooperates: embeddability in \texttt{M²} is exactly Dushnik–Miller dimension ≤ 2, recognizable in polynomial time [5, 6]. Above that the wall is real: the realizer route collapses (spacelike configurations of \texttt{M\textasciicircum{}{1+2}} induce circle containment orders, for whose recognition no polynomial algorithm is known and which are believed hard [7]), and the field's practical instruments — dimension estimators [8, 9], action-based tests [10], embedding procedures [11] — are either necessary conditions only, or procedures without theorems.

We therefore stop attacking symmetric decision and split the problem NP-style. A \textbf{certificate} of \texttt{d}-dimensional volume-faithfulness is a witness structure \texttt{W} such that \emph{(T1)} verifying \texttt{W} is polynomial-time and yields a constructive embedding with a defect bound, and \emph{(T2)} every faithful order possesses a \texttt{W} (its true geometric structure). Finding \texttt{W} is then honest heuristic territory: worst-case hard, as it must be — but \emph{any} success is a proof, not a score. Sections 2–3 build the certificate; Section 4 scopes its soundness exactly; Section 5 attacks the search face; Section 6 runs the \texttt{3+1} stack. Throughout, \texttt{D*} denotes anchored-box (star) discrepancy of a planar point set, and all gates, floors, and tolerances are calibrated and disclosed in the shipped receipts.

\section*{2. The two-dimensional engine}
Everything downstream scores 2D projections, so the 2D case must be solved \emph{per order}, deterministically and quantitatively. Prior results in this regime are typicality statements — the uniform measure on 2D orders concentrates on continuum-like orders [12], and random 2D orders are sprinklings of the diamond in distribution [13] — or embedding procedures without canonicity. We prove:

\textbf{Theorem 2.1 (canonical realizer).} A 2D order whose incomparability graph is prime (in the modular-decomposition sense) has exactly two transitive orientations (Gallai [14]), hence exactly one realizer pair \texttt{(L₁, L₂)} up to the parity swap \texttt{u ↔ v}. For non-prime orders the residual freedom is exactly the module freedom, and the \textbf{relabeling lemma} disposes of almost all of it: series modules admit no freedom, and a parallel block's quotient enters \texttt{L₁} in one order and \texttt{L₂} reversed, so reorienting it permutes which element carries which coordinate pair — the embedded point \emph{set}, hence \texttt{D*}, is exactly invariant. Only \emph{internal prime-node} flips can move the point set. \emph{(Demonstrated: bit-level invariance on a population carrying \texttt{105} modules at \texttt{N = 396}; the internal-prime residue closed by census — zero point-moving modules on sprinkling and lattice populations; on the one population where the freedom is real, \texttt{143} point-moving blocks are exhibited and priced at worst \texttt{|ΔD*| < 10⁻⁵}. Receipts \texttt{i1}, \texttt{i3}.)}

\textbf{Theorem 2.2 (volume gauge).} A 2D order's embedding freedom is exactly a pair of monotone reparametrizations of its two null coordinates — the discrete conformal freedom. Demanding uniform marginals fixes both: the \textbf{canonical rank embedding} (each null coordinate replaced by its normalized rank) is the volume-normalized representative of the order's conformal class, realizing the order exactly with exactly uniform marginals.

\textbf{Theorem 2.3 (characterization).} For a 2D order \texttt{P} on \texttt{n} elements, both directions quantitative: \emph{(forward)} if \texttt{P} admits a volume-faithful embedding with anchored-box defect \texttt{≤ ε}, then \texttt{dim(P) ≤ 2} and the canonical rank embedding has star discrepancy \texttt{D* ≤ 3ε + 1/n}; \emph{(backward, exact)} if the canonical rank embedding has \texttt{D* = δ}, then \texttt{P} is \texttt{δ}-volume-faithful — the rank embedding itself is such an embedding, and its defect \emph{is} its \texttt{D*}. Forward proof: any ε-faithful embedding's coordinate orders are the canonical realizer up to parity and relabeling-invariant module freedom, and monotone reparametrization to uniform marginals moves anchored-box mass by ≤ 2ε, plus the \texttt{1/n} rank discretization.

The theorem is constructive and the pipeline is polynomial (dimension-2 recognition and modular decomposition are classical polynomial problems [6, 15]; exact \texttt{D*} is \texttt{O(n² log n)} here). The showcase measurement: monotone distortions \texttt{(u², √v)} of a faithful embedding leave the order bit-identical while inflating the raw point-set discrepancy \texttt{0.053 → 0.301}; the canonical rank embedding recovers the volume gauge \emph{from the order alone} at \texttt{D* = 0.036} (receipt \texttt{i2}). Two honest boundaries, both measured: a bounded blow-up adversary (clusters of \texttt{c} co-located points) is flagged at finite \texttt{N} as a \texttt{√c} \texttt{D*}-inflation \emph{diagnostic} (measured \texttt{2.45×} against the predicted \texttt{√6 = 2.45}) but is asymptotically volume-faithful — its outright rejection belongs to a separate \emph{statistical} layer (interval-count dispersion; measured kill at \texttt{z = 28}); and volume-faithful ≠ Poisson-typical — a jittered lattice rightly passes the geometric characterization and fails the statistical layer (\texttt{z = 5.2} sub-ensemble dispersion). The characterization is a theorem about geometry; Poisson typicality is a model question deliberately split off.

\section*{3. Causal tomography in 2+1}
\textbf{Theorem 3.1 (shadows).} For \texttt{P} the order of a finite point set in \texttt{M\textasciicircum{}{1+(d−1)}}: for every spatial unit vector \texttt{m}, the projection \texttt{x ↦ (t, m·x⃗)} maps \texttt{P} to a 2D dominance order that coarsens \texttt{P}, with null coordinates the functionals \texttt{u\_{±m} = t ∓ m·x⃗}; and \texttt{P} is exactly the intersection of its shadow orders over the continuum celestial sphere (per pair, the supremum is achieved at \texttt{m = Δx⃗/|Δx⃗|}). \emph{(The coarsening half is exact in-receipt; the intersection identity's finite-\texttt{K} face is measured — a \texttt{K}-direction grid intersection converges to \texttt{P} from above, extra relations \texttt{32/11/1} of \texttt{17,291} comparable pairs at \texttt{K = 16/32/64}. Receipt \texttt{k1}.)}

A faithful diamond sprinkling's shadow is not uniform: its density is the \textbf{chord law} \texttt{ρ(t, s) ∝ √((1−|t|)² − s²)} (verified at total variation \texttt{0.023} against a 200k-point projection). The Section 2 certifier generalizes in one move: score each shadow's point set against the \emph{target copula} (the chord law pushed through its own quantile transforms — which also absorbs any monotone gauge freedom per direction) with a weighted star discrepancy \texttt{wD*}, floor-calibrated by same-law subsamples. Measured at \texttt{N = 400}: true shadows sit at \texttt{1.40×} the two-sample floor (gate \texttt{1.5×}), random-extension garbage at \texttt{9.4×}; the fine print is printed — the chord copula's departure from uniformity is small (\texttt{0.040} asymptotically), so at this \texttt{N} the weighting earns its keep in defect attribution and at scale, not as a dramatic finite-\texttt{N} kill.

\textbf{The harmonic condition.} For true Minkowski coordinates, \texttt{u\_θ(x) = t − x cos θ − y sin θ}: the direction profile of every point is \emph{exactly} a degree-≤1 Fourier polynomial in \texttt{θ}, and the coefficients \texttt{(a₀, −a₁, −b₁)} are \texttt{(t, x, y)}. The certificate condition demands quantile-gauged profiles (each \texttt{u\_θ} recovered intrinsically from ranks through the target marginal) be degree-≤1 concentrated. Measured: RMS residual \texttt{0.020}; coordinate reconstruction at correlation \texttt{0.9994} per axis; the rebuilt point set reproduces the original order on \texttt{99.67\%} of pairs and matches anchored-box counts to a mean absolute defect of \texttt{2.1} elements. Faithfulness transfers to the reconstruction with no closed-form volume formula anywhere.

\textbf{The celestial circle from inversion counts.} Classical multidimensional scaling on pairwise Kendall (inversion) distances between the \texttt{16} null-direction extensions recovers \texttt{S¹} — radius CV \texttt{0.009}, cyclic adjacency \texttt{16/16} — from order data alone, \emph{given the extension family}. Finding the family intrinsically is the search face (Section 5).

\textbf{Two measured honesty points.} The certifier has a \emph{plateau}: 400 adjacent-swap corruptions of the extensions leave the score at the floor, and corrupting all 16 directions still reconstructs coordinates at \texttt{> 0.999} per axis — certificate-grade families are interchangeable for reconstruction, which is aligned with the certificate's semantics (any certified witness yields the geometry). And annealing from a random start stalls at \texttt{6.0×} the floor after 12k local moves — the hardness face, exactly where the reframe put it. \emph{(Receipts \texttt{k2}, \texttt{k3} — the latter also demonstrates the chain-fleet transverse principle of Section 5: greedy chain extraction, cross-chain comparability coupling, MDS recovers the transverse plane at Procrustes \texttt{0.993}, shuffled-coupling control \texttt{0.19}.)}

\section*{4. Certificate v2 and the ghost budget}
\textbf{Certificate v2.} For a \texttt{d}-dimensional claim over direction set \texttt{{m\_k}}: \textbf{(i)} every antipodal shadow scores within a max-calibrated same-law floor (gate \texttt{1.3×} the floor in \texttt{2+1}, \texttt{1.35×} in \texttt{3+1}, replicate-calibrated); \textbf{(ii)} harmonic residual \texttt{< 0.05}; \textbf{(iii)} the harmonic reconstruction \emph{realizes the input order} at pair agreement \texttt{≥ 0.98}. Verdict bands, fixed in advance: \texttt{< 1.3×} certified; \texttt{1.3–2×} inconclusive; \texttt{> 2×} refused. Condition (iii) is what makes soundness discussable — the certificate asserts \emph{this order embeds as this point set}, not merely \emph{some projections look right} — and it is measurably load-bearing: a \textbf{separator adversary} (an extension family lifted from a \emph{different} faithful configuration) passes (i) at \texttt{0.040} and (ii) at \texttt{0.012} and fails (iii) at agreement \texttt{0.794} (receipt \texttt{l1}).

\textbf{The soundness budget, closed into a loop.} Discrete tomography's classical ghosts — switching components [16, 17] — dichotomize here: small-support rearrangements leave the configuration faithful up to negligible defect (the certificate rightly still passes — the input is simply another faithful order), while extended switching patterns \emph{are} density modulations. So the true adversary is spectral: \texttt{1 + a·cos(q⃗·x⃗)}. Two elementary facts pin its fate. \emph{(a)} Its shadow marginal in direction \texttt{n\_k} is fiber-averaged, leaving relative leakage \texttt{∼ a/(q·sin∠(q̂, n\_k)·ℓ)} — so invisibility to all \texttt{K} directions forces \texttt{q ≳ 1/sin(π/2K)}: Logan's finite-angle uncertainty principle [18] in our setting. \emph{(b)} Any mode's contribution to any anchored-box count — in a shadow or in the bulk — is at most \texttt{∼a/q}, because the CDF of an oscillation is flat to \texttt{1/q}. Together: \textbf{certificate-invisible ghosts cannot move the discrepancy metric anywhere}; what they can move is spectral density below wavelength \texttt{2π/q}. Demonstrated on all four legs (receipt \texttt{l1}): a \texttt{q = 120} ghost at \texttt{a = 0.6} is genuinely present (matched filter \texttt{z = 15.3}); it passes the full v2 certificate at \texttt{K = 8} \emph{and} its shadows at \texttt{K = 16} (more directions do not catch it, exactly as \emph{(b)} predicts) while bulk box discrepancy stays within its disclosed gate (\texttt{0.027 → 0.039}); and the same amplitude at \texttt{q = 6} is caught at over \texttt{2×} the floor. \textbf{The guarantee, final form:} certificate v2 implies an order-realizing embedding with anchored-box volume defect controlled at all scales; spectral non-uniformity above frequency \texttt{∼K} \emph{can be} invisible whenever its amplitude sits below the CDF floor — a budget that is metric-intrinsic (it binds any discrepancy-type certificate), amplitude-conditioned, and extendable by an explicit spectral axis (matched filters per frequency band) for callers who need it.

\section*{5. The intrinsic finder}
The search face closes to near-certificate range with a fully order-intrinsic pipeline, no annealing: \textbf{(1)} temporal seed \texttt{t̂ = |past(x)|\textasciicircum{}{1/d} − |future(x)|\textasciicircum{}{1/d}} with \texttt{d} the spacetime dimension (cone-volume counting: \texttt{1/3} in \texttt{2+1}, \texttt{1/4} in \texttt{3+1}); \textbf{(2)} transverse seed from a greedy chain fleet — extract long chains, form the cross-chain comparability coupling \texttt{w(c, c′)}, embed chains by MDS, extend to all elements by coupling-weighted averaging; \textbf{(3)} an order-only scale calibration (the MDS scale is arbitrary against \texttt{t̂}; a 12-point geometric scan on the certificate score fixes it — omitting this step scores \texttt{2.5×}, receipted); \textbf{(4)} the \textbf{tomographic bootstrap}: quantile-gauge each direction's profile through the target marginal, harmonically fuse into refined coordinates, iterate two to three times — the reconstruction operator as its own polish.

Measured in \texttt{2+1} (receipts \texttt{l2}, \texttt{m1}): at \texttt{N = 400}, \texttt{max wD* = 0.048} vs floor \texttt{0.033} — \texttt{1.5×}, inconclusive band, with condition (iii) at \texttt{0.953}; at \texttt{N = 1000} (\texttt{K = 12}), \texttt{0.0228} vs \texttt{0.0188} — \textbf{\texttt{1.21×}, inside the \texttt{1.3×} gate for condition (i)} — with (iii) at \texttt{0.9648}, still short of its \texttt{0.98} gate, so the overall verdict remains inconclusive and the residual is (iii)-only at this scale. The \texttt{1.5× → 1.21×} move measures the condition-(i) gap as finite-size, not structural. The refusal twin holds: the same pipeline refuses a Kleitman–Rothschild-type layered order (an explicit degenerate-fleet refusal path — such orders have no long-chain structure to seed from) and scores the cluster adversary at \texttt{2.2×} (refused band).

\section*{6. Three-plus-one dimensions}
The stack transfers with no structural change (receipt \texttt{l3}; \texttt{N = 900}, \texttt{12} Fibonacci-spiral directions plus antipodes). The shadow target becomes the \textbf{disk-section law} \texttt{(1−|t|)² − s²} (TV \texttt{0.013}); all \texttt{12} true shadows pass at the max-calibrated floor (gate \texttt{1.35×}, 10-replicate calibration, pass margin \texttt{+0.10} floors — small, and said so; fairly quantile-mapped garbage scores \texttt{4.2×} above the true band); the harmonic condition becomes degree-≤1 \emph{spherical} harmonics — profiles affine in \texttt{n ∈ S²}, residual \texttt{0.012} — whose four coefficients reconstruct \texttt{(t, x, y, z)} at \texttt{0.9999/0.9999/0.9998/0.9999} with the rebuilt order agreeing on \texttt{99.88\%} of pairs. The flagship: MDS on pairwise inversion distances between the \texttt{24} signed-direction extensions yields a 3D shell of radius CV \texttt{0.008}, a decisively 3-dimensional spectrum (fourth-to-third eigenvalue ratio \texttt{0.019}, gated on \texttt{|λ₄|/λ₃}), and pairwise angles matching the true celestial angles at correlation \texttt{0.9995} — to our knowledge the first reconstruction of the celestial sphere of a causal set from inversion counts. These families are \textbf{oracle-derived} (built from the true coordinates, as a verification of the certificate, not a search result), and we say so.

The intrinsic pipeline of Section 5 then runs end-to-end in \texttt{3+1} on the same configuration (receipt \texttt{m1}): seed quality \texttt{t̂} correlation \texttt{0.961} and transverse Procrustes \texttt{0.888} from \texttt{M = 26} chains (the 4D chains are shorter and the transverse space larger — the seed bars are deliberately lower and disclosed); after calibration and bootstrap, \texttt{max wD* = 0.0372} vs floor \texttt{0.0247} — \textbf{\texttt{1.50×}, condition (iii) \texttt{0.970}: inconclusive, at exact parity with the \texttt{2+1} pilot's own first stage}. Held-out reconstruction: \texttt{t} correlation \texttt{0.974}; spatial correlation \texttt{0.916} after Procrustes alignment (the recovered spatial frame is defined only up to global rotation, so per-axis spatial correlations are undefined under the certificate's own symmetry). Given Section 5's scale lesson, the larger-\texttt{N} \texttt{3+1} run is a compute question, not a structural one — but the certified end-to-end \texttt{3+1} claim is \textbf{not made here}.

\section*{7. Limitations and open problems}
All demonstrations are flat-diamond sprinklings at \texttt{N ≤ 1200}; curved backgrounds, boundaries, and non-diamond regions are untouched. The ghost budget's two legs are proved at the physics-paper level (fiber-averaging leakage; CDF flatness) with Logan cited for the principle; sharp constants and the finite-\texttt{N} noise-vs-leakage interplay are measured regimes, not theorems. The finder's refusal behavior is measured on two adversary classes, not proved. T1/T2 as stated in Section 1 are demonstrated constructively at pilot scale; their journal-grade quantitative forms (defect-propagation constants, the \texttt{S²} version of the harmonic argument) are open. The named open problems, in order of expected yield: (1) the certified end-to-end \texttt{3+1} reconstruction — close condition (iii)'s gap (\texttt{0.965/0.970} vs \texttt{0.98}) and run the \texttt{3+1} finder at larger \texttt{N}; (2) the sharp ghost-budget constants; (3) the statistical (Poisson-typicality) layer as a characterization rather than a battery of tests; (4) curved/conformally-flat extensions of the 2D engine.

\section*{8. Code and reproducibility}
All code is in \texttt{code/} and all canonical outputs in \texttt{receipts/} of this repository. Every number quoted above appears in a receipt output. The scripts are self-contained (NumPy only), run in minutes on a laptop, and print explicit \texttt{[PASS]/[FAIL]} gates with their tolerances; a script exits nonzero if any check fails.

\begin{itemize}
\item \texttt{i1\_canonical\_realizer.py} — Theorem 2.1: realizer canonicity, the relabeling lemma's bit-level invariance exhibit, Gallai verified live by small-\texttt{n} enumeration. Receipt: \texttt{i1\_canonical\_realizer.out}.
\item \texttt{i2\_characterization.py} — Theorems 2.2–2.3: the rank embedding, the \texttt{(u², √v)} distortion showcase (\texttt{0.301 → 0.036}), the blow-up diagnostic (\texttt{2.45×} vs \texttt{√6}), the statistical-layer split (lattice, \texttt{z = 5.2}). Receipt: \texttt{i2\_characterization.out}.
\item \texttt{i3\_internal\_prime\_census.py} — the internal-prime census closing Theorem 2.1's residue: absence on sprinkling/lattice populations; \texttt{143} point-moving blocks exhibited and priced \texttt{< 10⁻⁵} on the cluster population. Receipt: \texttt{i3\_internal\_prime\_census.out}.
\item \texttt{k1\_shadow\_engine.py} — Theorem 3.1's measured faces, the chord law (TV \texttt{0.023}), the weighted certifier with floor calibration, the harmonic reconstruction (\texttt{0.9994}/axis, \texttt{99.67\%} order agreement). Receipt: \texttt{k1\_shadow\_engine.out}.
\item \texttt{k2\_shadow\_search.py} — the celestial circle (CV \texttt{0.009}, adjacency \texttt{16/16}), the plateau and all-directions-corrupted reconstruction (\texttt{> 0.999}/axis), the cold-start stall (\texttt{6.0×}). Receipt: \texttt{k2\_shadow\_search.out}.
\item \texttt{k3\_decoration\_probe.py} — the chain-fleet transverse principle (Procrustes \texttt{0.993}; shuffled control \texttt{0.19}). Receipt: \texttt{k3\_decoration\_probe.out}.
\item \texttt{l1\_ghost\_budget.py} — certificate v2 jointly verified; the \texttt{q = 120} ghost's four legs; the \texttt{q = 6} control; the separator adversary (condition (iii) load-bearing). Receipt: \texttt{l1\_ghost\_budget.out}.
\item \texttt{l2\_intrinsic\_finder.py} — the Section 5 pipeline at \texttt{N = 400} (\texttt{1.5×}; naive-seed \texttt{2.5×}; refusals). Receipt: \texttt{l2\_intrinsic\_finder.out}.
\item \texttt{l3\_four\_dimensional.py} — the \texttt{3+1} stack (disk-section law; \texttt{12/12} shadows; spherical harmonics \texttt{0.9999}/axis; the celestial \texttt{S²}). Receipt: \texttt{l3\_four\_dimensional.out}.
\item \texttt{m1\_intrinsic\_4d.py} — the \texttt{2+1} scale test (\texttt{1.21×} at \texttt{N = 1000}) and the intrinsic \texttt{3+1} run (\texttt{1.50×}, (iii) \texttt{0.970}). Receipt: \texttt{m1\_intrinsic\_4d.out}.
\end{itemize}
\section*{References}
[1] L. Bombelli, J. Lee, D. Meyer, R. D. Sorkin, \emph{Space-time as a causal set}, Phys. Rev. Lett. \textbf{59}, 521 (1987).

[2] S. Surya, \emph{The causal set approach to quantum gravity}, Living Rev. Relativ. \textbf{22}, 5 (2019).

[3] L. Bombelli, J. Noldus, \emph{The moduli space of isometry classes of globally hyperbolic spacetimes}, Class. Quantum Grav. \textbf{21}, 4429 (2004).

[4] D. B. Malament, \emph{The class of continuous timelike curves determines the topology of spacetime}, J. Math. Phys. \textbf{18}, 1399 (1977).

[5] B. Dushnik, E. W. Miller, \emph{Partially ordered sets}, Amer. J. Math. \textbf{63}, 600 (1941).

[6] M. C. Golumbic, \emph{Algorithmic Graph Theory and Perfect Graphs} (Academic Press, 1980) — polynomial recognition of transitive orientability and dimension-2 orders.

[7] W. T. Trotter, \emph{Combinatorics and Partially Ordered Sets: Dimension Theory} (Johns Hopkins, 1992).

[8] J. Myrheim, \emph{Statistical geometry}, CERN preprint TH-2538 (1978).

[9] D. A. Meyer, \emph{The Dimension of Causal Sets}, PhD thesis, MIT (1988).

[10] D. M. T. Benincasa, F. Dowker, \emph{The scalar curvature of a causal set}, Phys. Rev. Lett. \textbf{104}, 181301 (2010).

[11] S. Johnston, \emph{Embedding causal sets into Minkowski spacetime}, arXiv:2111.09331; \emph{A new algorithm for embedding causal sets}, arXiv:2502.09701.

[12] G. Brightwell, J. Henson, S. Surya, \emph{A 2D model of causal set quantum gravity: the emergence of the continuum}, Class. Quantum Grav. \textbf{25}, 105025 (2008).

[13] P. Winkler, \emph{Random orders}, Order \textbf{1}, 317 (1985).

[14] T. Gallai, \emph{Transitiv orientierbare Graphen}, Acta Math. Acad. Sci. Hungar. \textbf{18}, 25 (1967).

[15] R. M. McConnell, J. P. Spinrad, \emph{Modular decomposition and transitive orientation}, Discrete Math. \textbf{201}, 189 (1999).

[16] H. J. Ryser, \emph{Combinatorial properties of matrices of zeros and ones}, Canad. J. Math. \textbf{9}, 371 (1957).

[17] R. J. Gardner, P. Gritzmann, \emph{Discrete tomography: determination of finite sets by X-rays}, Trans. Amer. Math. Soc. \textbf{349}, 2271 (1997).

[18] B. F. Logan, \emph{The uncertainty principle in reconstructing functions from projections}, Duke Math. J. \textbf{42}, 661 (1975). (Both the ghost budget here and the longest-chain constants of causal-set dimension estimation trace to B. F. Logan — see also B. F. Logan, L. A. Shepp, Adv. Math. \textbf{26}, 206 (1977).)

\end{document}
